\begin{lemma}
\label{lemma-intersection-subfields}
Let $K$ be a field of characteristic $p$. Let $\{K_\alpha\}_{\alpha \in A}$
be a collection of subfields of $K$ with the following properties
\begin{enumerate}
\item $K^p \subset K_\alpha$ for all $\alpha \in A$,
\item $K^p = \bigcap_{\alpha \in A} K_\alpha$,
\item for $\alpha, \alpha' \in A$ there exists an $\alpha'' \in A$
such that $K_{\alpha''} \subset K_\alpha \cap K_{\alpha'}$.
\end{enumerate}
Then
\begin{enumerate}
\item the intersection of the kernels of the maps
$\Omega_{K/\mathbf{F}_p} \to \Omega_{K/K_\alpha}$ is zero,
\item for any finite extension $L/K$ we have
$L^p = \bigcap_{\alpha \in A} L^pK_\alpha$.
\end{enumerate}
\end{lemma}

\begin{proof}
Proof of (1).
Choose a $p$-basis $\{x_i\}$ for $K$ over $\mathbf{F}_p$.
Suppose that $\eta = \sum_{i \in I'} y_i \text{d}x_i$ maps to zero in
$\Omega_{K/K_\alpha}$ for every $\alpha \in A$. Here the index set
$I'$ is finite. By Lemma \ref{lemma-p-basis}
this means that for every $\alpha$ there exists a relation
$$
\sum\nolimits_E a_{E, \alpha} x^E = 0,\quad a_{E, \alpha} \in K_\alpha
$$
where $E$ runs over multi-indices $E = (e_i)_{i \in I'}$ with
$0 \leq e_i < p$. On the other hand, Lemma \ref{lemma-p-basis}
guarantees there is no such relation $\sum a_E x^E = 0$ with
$a_E \in K^p$. This is a contradiction by
Lemma \ref{lemma-intersection-subfields-subspace}.

\medskip\noindent
Proof of (2). Suppose that we have a tower $L/M/K$
of finite extensions of fields. Set $M_\alpha = M^p K_\alpha$
and $L_\alpha = L^p K_\alpha = L^p M_\alpha$. Then we can first prove that
$M^p = \bigcap_{\alpha \in A} M_\alpha$, and after that prove
that $L^p = \bigcap_{\alpha \in A} L_\alpha$. Hence it suffices to
prove (2) for primitive field extensions having no nontrivial subfields.
First, assume that $L = K(\theta)$ is separable over $K$. Then
$L$ is generated by $\theta^p$ over $K$, hence we may assume that
$\theta \in L^p$. In this case we see that
$$
L^p = K^p \oplus K^p\theta \oplus \ldots K^p\theta^{d - 1}
\quad\text{and}\quad
L^pK_\alpha =
K_\alpha \oplus K_\alpha \theta \oplus \ldots K_\alpha\theta^{d - 1}
$$
where $d = [L : K]$. Thus the conclusion is clear in this case. The other
case is where $L = K(\theta)$ with $\theta^p = t \in K$, $t \not \in K^p$.
In this case we have
$$
L^p = K^p \oplus K^pt \oplus \ldots K^pt^{p - 1}
\quad\text{and}\quad
L^pK_\alpha =
K_\alpha \oplus K_\alpha t \oplus \ldots K_\alpha t^{p - 1}
$$
Again the result is clear.
\end{proof}